\defmodule{palpha}

\baselineskip=14pt

Le $P_\alpha(s)$ est une mesure de la qualit\'e d'int\'egration d'une
r\`egle de treillis.  D'une mani\`ere g\'en\'erale, 
une r\`egle de treillis permet d'approximer
l'int\'egrale de $f$, une fonction, sur le cube unitaire de dimension $s$.
$f$ doit \^etre lisse et p\'eriodique par rapport \`a $x \in$ {\bf R}.

La formule g\'en\'erale du $P_\alpha(s)$ pour $\alpha$ entier et pair est 
\eq
  P_\alpha (s) = -1 + 1/m \sum_{i=0}^{m-2} \prod_{j=0}^{s-1} \Big[
                     1- \frac{(-1)^{\alpha /2}(2\pi)^{\alpha}}{\alpha !}B_\alpha(x_{i+j}) \Big] 
										\label{eq:Palpha}
\endeq
o\`u les $B_\alpha(x)$ sont les polyn\^omes de Bernoulli. $x_{i+j}$ est le $i+j^e$ nombre
g\'en\'er\'e et se situe entre 0 et 1.
Les premiers polyn\^omes de Bernoulli sont \\\\ 
\hspace*{18mm} \baselineskip=18pt
$B_0(x)=1$,\hspace*{9mm} $B_1(x)=x-1/2$,\hspace*{9mm}
$B_2(x)=x^3-x^2+1/6$,  \\ \hspace*{18mm}
$B_3(x)=x^3-(3/2)x^2+(1/2)x$,\hspace*{5mm} $B_4(x)=x^4-2x^2+x-1/30$

\noindent
Les $B_\alpha(x)$ se calculent en utilisant l'extension de la s\'erie suivante:
$$\frac{te^{xt}}{e^t-1} = \sum_{s=0}^{\infty}B_s \frac{t^s}{s!}$$
On peut ajouter un poids \`a chacune des dimensions $s$, pour des $\alpha$ pairs.
La formule du
$P_\alpha(s)$ devient: 
$$P_{2\alpha}(s) = \beta_0 \Big\{ -1 + 1/m \sum_{i=0}^{m-2} \prod_{j=0}^{s-1} \Big[
                     1- \frac{\big( -(2\pi\beta_{j+1})^2\big) ^{\alpha}}{2\alpha !}B_{2\alpha}(x_{i+j}) \Big] \Big\} 
$$

Le $B_\alpha(s)$, \`a ne pas confondre avec le $B_\alpha(x)$, est une mesure
indiquant qu'il existe au moins un $P_\alpha(s) \le B_\alpha(s)$.  D'une mani\`ere
g\'en\'erale, on peut le voir comme une borne sup\'erieure valide en pr\'esence
des conditions suivantes: $m> e^{\alpha s/(\alpha -1)}$ et premier, $s\ge 2$ et
$\alpha > 1$.  Pour plus de d\'etails consulter le th\'eor\`eme 4.4, p.83 du
livre de Sloan et Joe. 
Il se calcule ainsi:
\begin{eqnarray*}
B_\alpha(s) = {\Big( \frac{e}{s} \Big)}^{\alpha s} \frac{(2\ln m + s)^{\alpha s}}{m^\alpha}
\end{eqnarray*} 

\large
{\bf Le module PALPHA}
\normalsize

Module servant de support aux programmes PAl, BAl et Coeff. Ils permettent
respectivement de trouver le $P_\alpha$ ou  le $B_\alpha$ d'un g\'en\'erateur
ou enfin de trouver des bons coefficients pour un g\'en\'erateur en
utilisant le $P_\alpha$ comme discriminant.  

\large
{\bf palpha}
\normalsize

Ce programme permet de calculer le $P_\alpha$ pour un 
g\'en\'erateur et un $\alpha$ donn\'e. 
Tout d\'ependant du
$\alpha$
la fonction de calcul changera. Les donn\'ees suivantes doivent \^etre
plac\'ees dans un fichier avec l'extension {\tt".dat"} selon le format de
la figure~1.  Les r\'esultats seront plac\'es dans un fichier  avec l'extension {\tt".res"}.

\large
{\bf balpha}
\normalsize

Permet de calculer le $B_\alpha$ d'un g\'en\'erateur.  Pour l'instant
seulement le cas avec $\alpha = 2$ est impl\'ement\'e.  Le format du
fichier est le m\^eme que pour le programme PAl. Les donn\'ees devront
\^etre plac\'ees dans un fichier avec l'extension {\tt".dat"} et les
r\'esultats se trouveront dans le fichier avec le m\^eme pr\'efixe,
cette fois avec l'extension {\tt".res"}.

\noindent
{\bf Le format du fichier d'entr\'ee pour PAl ou BAl}

\begin{center}
\it{
\begin{tabbing}
\hspace{2cm}\= LCG \hspace{3cm} \= SorteGenerateur \\
\>PAl ou BAl \>   Type de calcul \\
\>251     \>      m  modulo\\
\>true    \>      modulo premier\\
\>33      \>      a  multiplicateur\\
\>2       \>      mindim\\
\>4      \>      maxdim\\
\>2      \>      saut \\ 
\>2       \>      alpha\\
\>false   \>      VerifBooleen\\
\>1       \>      seed\\
\>1.0 2.0 3.0 1.0 1.0 \> Beta(0),...,Beta(MaxDim)\\
\end{tabbing}
}
\rm {Figure 1 : Format du fichier pour PAl ou BAl}
\end{center}

%%\begin{figure}[h]
%%\caption{Format du fichier pour PAl ou BAl}
%%\end{figure}

\noindent
{\it Type de calcul} : Permet de d\'eterminer le type de calcul
\`a effectuer.

\noindent
{\it Modulo} : le modulo du g\'en\'erateur en base 10, soit un nombre premier ou
une puissance de 2.  Ce nombre doit \^etre inf\'erieur \`a 3.03E09 ou $2^{31}$.

\noindent
{\it Modulo premier} : Bool\'een indiquant si le modulo utilis\'e est un 
nombre premier.

\noindent
{\it Multiplicateur} : le multiplicateur du g\'en\'erateur en base 10. Il doit \^etre
inf\'erieur au modulo.  Ou un intervalle dans le cas de la recherche de coefficient.
La recherche permettra de trouver le meilleur coefficient dans cet intervalle en fonction du
$P_\alpha$.

\noindent
{\it Mindim} : Entier  donnant la borne inf\'erieur sur le s (la dimension) 
dans le calcul \`a effectuer.  

\noindent
{\it Maxdim} : Entier  donnant la borne sup\'erieur sur le s (la dimension)
dans le calcul \`a effectuer.

\noindent
{\it Saut} : Entier d\'eterminant le saut \`a effectuer entre les dimensions.

\noindent
{\it alpha} : Entier pair, en g\'en\'eral de 2 \`a 10.

\noindent
{\it VerifBooleen} : Si {\tt true} alors on v\'erifira si la p\'eriode est maximale,
mettre \`a {\tt false} si aucune v\'erification n'est n\'ecessaire.  Dans le cas de 
la recherche de coefficient, si {\tt true} alors seulement les g\'en\'erateurs
\`a p\'eriode maximale seront consid\'er\'es, si {\tt false} alors le
$P_\alpha$ de tous les g\'en\'erateurs sera \'evalu\'e.  Si par contre le
modulo n'est pas un nombre premier alors peu importe la valeur entr\'ee \`a
{\it VerifBooleen},
nous allons effectuer une recherche sur l'ensemble des coefficients.


\noindent
{\it Seed} : Entier inf\'erieur au modulo qui servira de germe dans le g\'en\'erateur.  

\noindent
{\it Beta} : Nombres r\'eels permettant de donner un poids \`a chacune des dimensions.
Il faut entrer dans le fichier une valeur pour $\beta (0),...,\beta (MaxDim)$. 
